% ======================================================================
\section{Concluding remarks}\label{sec:discussion}
% ======================================================================

\subsection{The disordered side}

Three quantities determine the lower bound: the vertex gain $\gamma$, the chord price $\theta$ and the growth rate of the shape-weighted path count. Each can be improved independently.

The gain $\gamma$ is set by the prism, which treats every site outside its ten specified sites as occupied. A larger local graph, or one that records the occupation of second neighbours, increases the certified local integrals. Joint gains from pairs of nearby middle bonds are a further source. At $r=309/200=1.545$ even the most favourable degree-two integrals give a shape-weighted growth factor $u\rho>1.0039$ on the $16$-step automaton. Gains of this kind are therefore what is needed to move the lower bound substantially further.

The chord price enters through the H\"older step of \cref{lem:holder}, which converts the chord moment, growing like $5^n$, into the factor $e^{\theta/4}$. A transfer operator that records contacts and shapes simultaneously avoids this conversion. The contact operator of \cref{prop:r1} already records contacts. At $r=1.545$, contact operators with $w=4/5$ and $w=6/5$ show that paths with chord density $c(P)/|P|\le0.137$, and those with $c(P)/|P|\ge0.253$, give a uniformly summable contribution to the susceptibility.

Finally, Fisher's bound evaluated at the numerical value $\mu=4.684039931(27)$ of the connective constant~\cite{Cli13} would give $r_c>1.5428$. Combining the factor $e^{-I(P,u)}$ with sharper walk counting is therefore a natural route to further improvement.

\subsection{The ordered side}

The upper bound rests on a linear programme. Its constraints are exact identities and inequalities for local correlations, and its unknown is the zero mode $q$. The infrared bound of \cref{prop:infrared} enters through five test polynomials. The law of the shell spins enters through the mixture representation of \cref{lem:mixture}, with constraints that couple field labels across single links and corner sources through single flips. Additional valid local inequalities, a larger shell, or constraints coupling several labels at once would each enlarge the ordered region certified by the same method.
